\section{Introduction}

Many sequential decision-making problems require solving an optimization problem whose parameters are not known in advance but must be predicted from contextual data. The standard machine learning pipeline first trains a predictive model to minimize forecast error, and then feeds the predictions into a downstream optimizer. However, minimizing prediction error does not, in general, minimize the cost of the decisions that result from those predictions~\cite{}. Recently, decision-focused learning (DFL)~\cite{} has been proposed to address this mismatch by training the predictive model end-to-end: gradients of the downstream decision quality are backpropagated through the optimization layer to update the model parameters directly.

When the downstream problem is a Markov decision process (MDP), the optimization layer solves for an optimal policy given predicted rewards or transition probabilities. Existing DFL methods for MDPs typically differentiate through the Karush--Kuhn--Tucker (KKT) conditions of the Bellman optimality equations, treating the policy as the solution of a nonlinear system (cite Kai's paper). This requires solving a linear system whose dimension scales with $|\gS||\gA|$, the product of the state and action space sizes. For structured control problems with large action spaces, the resulting computational and numerical costs can be prohibitive.

In this work, we take a different route to exploit the \emph{linear programming} (\LP) formulation of MDPs over occupancy measures. Under the strictly positive initial distribution, the LP is guaranteed to be non-degenerate: its optimal basis is uniqure, and the variable dimension is $|\gS| \times |\gS|$, leading to significant computational reduction over KKT-based approaches, especially when the action space $|\gA|$ is large.

A second challenge in DFL for MDPs is that the decision landscape is piecewise linear: the optimal policy is constant within each basis region of the LP, and small perturbation to the predicted parameters can trigger a basis change. As a result, its gradient is locally constant almost everywhere and discontinuous at the basis boundaries, a known source of zero or uninformative gradient in DFL with linear objectives~\citep{mandi2025minimizing, berthet2020learning, elmachtoub2022smart}. We address this by introducing \textbf{state aggregation} into the DFL framework. 

Finally, we evaluate that our proposed methods on multiple tasks and show it outperforms all the baselines.

Contributions:
1. A novel LP occupancy measure formulation for DFL gradients in a finite discounted MDP that leads to a very efficient derivative computation.
2. State aggregation that makes the loss landscape smoother with better regrets and some theoretical guarantees.
3. Multiple real-world applications.